% sections/method.tex

\section{Method}
\label{sec:method}

CARE-Judge converts an LLM judge $f$ into a \emph{selective evaluator} $(\hat{f}, \hat{c})$ that either outputs a judgment $\hat{f}(x)$ or abstains, based on a confidence measure $\hat{c}(x)$ and a risk-controlled threshold $\hat{\lambda}$.
The framework has three components: (1)~multi-signal protocol-stability feature extraction, (2)~learned correctness calibration, and (3)~risk-controlled threshold selection with finite-sample guarantees.

\subsection{Protocol-Stability Features}
\label{sec:method:features}

Given a pairwise evaluation item $x = (q, a_1, a_2)$ consisting of a prompt $q$ and two candidate responses, we issue multiple judge calls under semantically equivalent evaluation protocols and measure the stability of the resulting verdicts.
Let $f(x; r, \tau)$ denote the judge's output (winner $\in\{A, B\}$ and confidence $\in[0,1]$) given rubric $r$ and sampling temperature $\tau$.

\paragraph{Rubric Perturbation Stability.}
We define a set of $K=3$ semantically equivalent rubric variants $\{r_1, \ldots, r_K\}$ that are intended to express the same evaluation criteria in different wording (e.g., ``Prefer the more correct response'' vs.\ ``Choose the answer a careful evaluator would prefer''); all variants are listed verbatim in the Appendix.
For each item, we collect verdicts $\{v_1, \ldots, v_K\}$ where $v_k = f(x; r_k, 0)$.
We compute:
\begin{align}
    \text{rubric\_vote\_share} &= \max_{w \in \{A,B\}} \frac{|\{k : v_k = w\}|}{K}, \label{eq:rubric_share} \\
    \text{rubric\_flip} &= \mathbb{1}\!\left[|\{v_1, \ldots, v_K\}| > 1\right], \label{eq:rubric_flip}
\end{align}
where $\text{rubric\_flip} = 1$ if any rubric variant produces a different winner.
The intuition is that a verdict that changes under paraphrase is \emph{protocol-dependent} and therefore not robust to the stated evaluation specification. We deliberately avoid the stronger claim that such a verdict is ``content-irrelevant,'' because a paraphrase can also shift the relative weight the judge places on competing criteria (helpfulness, clarity, style) rather than only its sensitivity to surface wording; the Appendix reports a negative-control experiment with intentionally non-equivalent rubrics to separate wording perturbation from criterion perturbation.

\paragraph{Position-Swap Consistency.}
We judge the item in both response orderings: $(q, a_1, a_2)$ and $(q, a_2, a_1)$.
Let $v_{\text{orig}} = f(x; r_1, 0)$ and $v_{\text{swap}} = f(\text{swap}(x); r_1, 0)$, where $\text{swap}$ reverses the response order.
We map $v_{\text{swap}}$ back to the original coordinate system and compute:
\begin{equation}
    \text{swap\_consistency} = \mathbb{1}\!\left[\text{map}(v_{\text{swap}}) = v_{\text{orig}}\right].
    \label{eq:swap}
\end{equation}
This probes positional bias: a judge that reverses its verdict under response reordering is unreliable.

\paragraph{Self-Consistency.}
We form a set of $S$ judgments under the base rubric that pools the deterministic base verdict ($\tau{=}0$) with $S-1$ stochastic re-samples at temperature $\tau>0$, and compute the majority vote share:
\begin{equation}
    \text{self\_vote\_share} = \max_{w} \frac{|\{s : f(x; r_1, \tau_s) = w\}|}{S}.
    \label{eq:self}
\end{equation}
In all reported experiments $S=3$ (one $\tau{=}0$ base verdict and two $\tau{=}0.7$ samples), so the score is genuinely informative rather than degenerate.
An optional \emph{adaptive-$k$} early-stopping rule ($\text{self\_vote\_share}\geq\tau_{\text{stop}}$ after a minimum of two samples) is available to reduce cost but was disabled in the main experiments.

\paragraph{Simulated Annotators (optional; disabled in main runs).}
Following~\citet{jung2025trust}, the framework can construct $N$ few-shot judge personas conditioned on labeled demonstrations and use their agreement as an additional signal. To avoid calibration leakage, any such demonstrations must be drawn only from the training split $D_{\text{train}}$, never from the calibration split $D_{\text{cal}}$ used for thresholding (Section~\ref{sec:method:risk}). Because this signal reproduces Trust-or-Escalate rather than testing our hypothesis, we disable it in all main experiments ($N{=}0$); the corresponding feature is then a constant fallback and carries no information. We include it only in a supplementary comparison against a simulated-annotator-only baseline.

\paragraph{Feature Vector.}
With simulated annotators disabled, CARE-Judge uses six raw judge calls per item---one base ($\tau{=}0$), two self-consistency samples ($\tau{=}0.7$), one position swap, and two additional rubric paraphrases (the base rubric doubling as the third rubric verdict)---and forms the feature vector
\begin{equation}
\begin{aligned}
    \boldsymbol{\phi}(x) = \big[\,
    &\text{base\_conf},\; \text{rubric\_vote\_share},\; \text{rubric\_flip},\\
    &\text{swap\_consistency},\; \text{self\_vote\_share},\; \text{sim\_vote\_share},\\
    &\text{mean\_conf},\; \text{std\_conf},\; \text{length\_gap}\,\big],
\end{aligned}
    \label{eq:features}
\end{equation}
where $\text{sim\_vote\_share}$ is the constant fallback noted above.

\subsection{Correctness Calibration}
\label{sec:method:calibration}

We fit a calibrator $\hat{p}: \mathbb{R}^d \to [0,1]$ that predicts the probability that the judge's verdict is correct given the feature vector:
\begin{equation}
    \hat{p}(x) = P\!\left(\hat{f}(x) = y \mid \boldsymbol{\phi}(x)\right),
    \label{eq:calibrator}
\end{equation}
where $y$ is the reference label.
We support logistic regression, isotonic regression, and gradient-boosted classifiers.
The calibrator is fit on a \emph{training split} $D_{\text{train}}$ that is disjoint from both the threshold-selection set and the evaluation set.

\subsection{Risk-Controlled Threshold Selection}
\label{sec:method:risk}

Given a user-specified risk tolerance $\alpha \in (0,1)$ and error level $\delta \in (0,1)$, we select an acceptance threshold $\hat{\lambda}$ on a \emph{calibration split} $D_{\text{cal}}$ (disjoint from $D_{\text{train}}$) such that:
\begin{equation}
    P\!\left(P\!\left(\hat{f}(x) \neq y \mid \hat{c}(x) \geq \hat{\lambda}\right) \leq \alpha\right) \geq 1 - \delta.
    \label{eq:guarantee}
\end{equation}

We employ fixed-sequence testing~\citep{bauer1991multiple,jung2025trust}: candidate thresholds are ordered a priori from most to least confident, and we walk down this pre-specified sequence, at each step testing whether the $(1-\delta)$ upper confidence bound on the accepted-set error rate is $\leq\alpha$. We accept the lowest threshold that still satisfies the bound and stop at the first subsequent violation, which controls the family-wise error without a separate multiplicity correction.
We use the \emph{exact} one-sided Clopper--Pearson upper bound~\citep{clopper1934use}:
\begin{equation}
    \text{CP}_{\text{upper}}(e, n, \delta) = \sup\!\left\{p : P\!\left(\text{Bin}(n, p) \leq e\right) \geq \delta\right\},
    \label{eq:cp}
\end{equation}
where $e$ is the number of errors among the $n$ accepted calibration items.
This bound is computed exactly via bisection on the binomial CDF.
Algorithm~\ref{alg:care} states the complete train/calibrate/test procedure, including the pre-specified stopping rule and the minimum-acceptance floor $m$ that prevents accepting a threshold on too few calibration points.

\begin{algorithm}[t]
\caption{CARE-Judge selective evaluation}
\label{alg:care}
\begin{algorithmic}[1]
\STATE \textbf{Input:} labeled items $D$; risk $\alpha$, level $\delta$; min-accept $m$
\STATE Split $D$ into disjoint $D_{\text{train}}, D_{\text{cal}}, D_{\text{test}}$ (40/30/30)
\STATE Fit calibrator $\hat p$ on $\{(\boldsymbol\phi(x), \mathbb{1}[\hat f(x){=}y])\}_{x\in D_{\text{train}}}$
\STATE Sort $D_{\text{cal}}$ by $\hat p$ descending: $c_{(1)}\ge\cdots\ge c_{(n)}$
\STATE $\hat\lambda \gets \infty$ \COMMENT{accept nothing by default}
\FOR{$i = 1$ \TO $n$}
  \STATE $e_i \gets \#\{j\le i: \hat f(x_{(j)})\ne y_{(j)}\}$
  \IF{$i \ge m$ \AND $\text{CP}_{\text{upper}}(e_i, i, \delta) \le \alpha$}
    \STATE $\hat\lambda \gets \hat p(x_{(i)})$ \COMMENT{lower threshold, more coverage}
  \ELSIF{$i \ge m$ \AND $\hat\lambda < \infty$}
    \STATE \textbf{break} \COMMENT{stop at first violation after acceptance}
  \ENDIF
\ENDFOR
\STATE \textbf{Output:} on $D_{\text{test}}$, accept $x$ iff $\hat p(x)\ge\hat\lambda$, else abstain
\end{algorithmic}
\end{algorithm}

The selective evaluator is then applied to a \emph{held-out test split} $D_{\text{test}}$:
\begin{equation}
    (\hat{f}, \hat{c})(x) = \begin{cases} \hat{f}(x) & \text{if } \hat{c}(x) \geq \hat{\lambda}, \\ \emptyset & \text{otherwise (abstain).} \end{cases}
    \label{eq:selective}
\end{equation}

All reported metrics (coverage, selective risk, accepted accuracy, calibration error) are measured on $D_{\text{test}}$, ensuring the risk guarantee is valid.

\subsection{Cost-Aware Cascade (Exploratory)}
\label{sec:method:cascade}

As an application, we arrange multiple judges $\mathcal{M} = (M_1, \ldots, M_T)$ in order of increasing cost. Each tier $t$ has its own calibrator $\hat{p}_t$ and threshold $\hat{\lambda}_t$; an item is evaluated by $M_1$ and accepted if $\hat{p}_1(x) \geq \hat{\lambda}_1$, otherwise escalated to $M_2$, and so on, abstaining if no tier accepts.
To make the cascade guarantee sound, we calibrate each tier \emph{conditionally}: tier $t$'s threshold is selected on the subset of $D_{\text{cal}}$ that was rejected by tiers $1{:}t{-}1$, i.e.\ on the distribution of items that actually reach $M_t$ at deployment. We also split the error budget as $\delta_t = \delta/T$, so that by a union bound the event ``some tier's risk bound fails'' has probability at most $\delta$, giving a simultaneous $(\alpha,\delta)$ guarantee across the whole cascade. We still treat the cascade as an exploratory application (Appendix~\ref{app:cascade}), because it uses a single labeled pool per dataset and we have not stress-tested the conditional calibration under distribution shift; a formal deployment-time theorem under shift is future work.
